\subsection{A full gentle element in the extra CM coset}

\begin{statement}{E.58}{Proved}{Relative pronilpotent lifting and the extra CM coset}
Let \(\Gamma\) be either the free profinite pants group
\(\widehat F(a,b)\), with peripherals \(a,b,(ab)^{-1}\),
or the actual modular group
\(\mathcal B=C_2\amalg C_3=\langle s,r\mid s^2=r^3=1\rangle\),
with peripherals \(s,r,(sr)^{-1}\).
Let \(U\triangleleft\Gamma\) be an open normal subgroup, and put
\[
 \begin{gathered}
 D=\Gamma/U,\qquad R_U=R_{\rm nil}(U),\qquad
 \Gamma_{U,\rm nil}=\Gamma/R_U,\qquad
 \Gamma_{U,p}=\Gamma/O^p(U).
 \end{gathered}
 \tag{E.58.1}
\]
Here all residuals are closed, and \(O^p(U)\) is the kernel
of the maximal pro-\(p\) quotient of the actual subgroup \(U\).
Every continuous positive peripheral automorphism of
\(\Gamma_{U,\rm nil}\), with any specified three conjugators
there, lifts to a continuous automorphism of \(\Gamma\) with
those full conjugator coordinates retained. A specified peripheral
whose old conjugator is one can be fixed literally in the lift.
The exact local-source comparison is
\[
 \Gamma_{U,\rm nil}
       \ \xrightarrow{\ \sim\ }\ 
       \mathop{\prod}_{p,\,D}\Gamma_{U,p}.
 \tag{E.58.2}
\]
The product on the right consists of tuples with one common value
in the specified finite quotient \(D\).

In the conventions of E.1, set
\[
 \begin{gathered}
 \tau=sr,\qquad F=\ker(\mathcal B\to S_3),\qquad
 x=\tau^{-2},\quad y=sxs,\qquad H=F'F^3,\\
 q:\mathcal B\twoheadrightarrow
 Q=\operatorname{SL}_2(\mathbf Z/12)/\{\pm I\},\\
 q(s)=\left[\begin{pmatrix}0&1\\-1&0\end{pmatrix}\right],\quad
 q(r)=\left[\begin{pmatrix}0&-1\\1&-1\end{pmatrix}\right].
 \end{gathered}
 \tag{E.58.3}
\]
There is an actual full profinite gentle pair \((1,f)\), with
\[
 f\in R_{\rm nil}(H)\subset F',
 \tag{E.58.4}
\]
whose normalized modular automorphism \(\Phi\) has actual
conjugators \(A,\beta,V\in F\) satisfying
\[
 \begin{gathered}
 \Phi(r)=ArA^{-1},\qquad
 \Phi(s)=\beta s\beta^{-1},\\
 \Phi(\tau)=V\tau V^{-1}=\tau,\qquad V=1,\\
 q(A)=[5I],\qquad q(\beta)=q(V)=1.
 \end{gathered}
 \tag{E.58.5}
\]
Its two E.1 CM bits are \((1,0)\). Thus
\(\widehat{\mathrm{GT}}_0\) reaches the extra CM coset,
and Task E.2 has a positive answer.
The construction is an actual full-source lift with these
specified finite frames.
\end{statement}

\begin{proof}
We first prove the relative source theorem, then construct its
particular centralizer coordinate and calculate the CM signs.

\paragraph{The complete relative lifting theorem.}
Because \(U\) is normal, its characteristic residual \(R_U\)
is normal in \(\Gamma\), so all quotients in
\textup{(E.58.1)} are actual continuous quotients of the same source.
Let \(\alpha\) be the specified peripheral automorphism of
\(\Gamma_{U,\rm nil}\), with conjugators \(b_i\).
Work in the compact space \(\Gamma^3\) of their prospective
full lifts. Impose the peripheral product relation and generation
in every finite quotient, and the values of the three conjugators
in every finite quotient of \(\Gamma_{U,\rm nil}\).

Take one actual finite joint source
\(q_0:\Gamma\twoheadrightarrow E_0\) containing any prescribed
finite collection of these constraints and the quotient \(D\).
Write
\[
 U_0=q_0(U)=\ker(E_0\to D),\qquad
 A_0=\gamma_\infty(U_0),\qquad E_1=E_0/A_0.
 \tag{E.58.6}
\]
The subgroup \(A_0\) is characteristic in \(U_0\), hence normal
in \(E_0\), and
\([E_0,A_0]=A_0\), because already \([U_0,A_0]=A_0\).
Apply the finite scalar-cover construction D.50.5--17 to this
actual extension. It gives an actual source quotient
\(E_1^\dagger\) central over \(E_1\), including every subgroup
count and every conjugator coordinate used there.
Its image of \(U\) is a central extension of the nilpotent group
\(U_0/A_0\), and is therefore nilpotent. Consequently its source
map kills \(R_U\):
\[
 \Gamma\longrightarrow\Gamma_{U,\rm nil}
                  \longrightarrow E_1^\dagger.
 \tag{E.58.7}
\]
This proves the required factorization for this particular
auxiliary quotient, including its original \(D\)-coordinate.

Precompose \textup{(E.58.7)} by \(\alpha\). The new source map
is onto and has the specified images of the \(b_i\) as literal
peripheral conjugators. D.50.15--17 then supplies an epimorphism
to \(E_0\), with those conjugator values retained in every old
factor. Every old relative finite constraint kills \(A_0\),
since the image of \(U\) in that factor is nilpotent.
Thus this one finite epimorphism meets the entire chosen finite
collection. In the modular case its source really is
\(C_2\amalg C_3\): D.50.9 constructs the auxiliary maps with
the two actual torsion relations, and the counted boundary
classes retain these relations.

Compactness now gives three full conjugators and an endomorphism
of \(\Gamma\) generating every finite quotient. Its image is
compact, hence is all of \(\Gamma\). The source is finitely
generated, so the endomorphism is an automorphism by the
Hopficity argument in D.50. This proves the full commuting
quotient square and every specified conjugator value.
If the conjugator of a retained peripheral is one, its full
lift lies in \(R_U\); conjugating the full automorphism by
its inverse fixes that peripheral literally and leaves all
the retained coordinates unchanged.
This proof applies in particular to characteristic \(U\);
it also proves the stated normal-subgroup version with its
specified finite base.

\paragraph{The exact product of the relative local sources.}
The actual map
\[
 U/R_{\rm nil}(U)\longrightarrow\prod_p U^{(p)}
 \tag{E.58.8}
\]
is an isomorphism. Its kernel is the nilpotent residual,
because a finite nilpotent group is the product of its Sylow
groups. For surjectivity, every finite product of finite
quotients at distinct primes is the actual source image:
Goursat's common quotients have coprime orders and are trivial.
The image in the infinite product is compact, so these finite
surjections give the full product.

The kernel of \textup{(E.58.2)} is therefore exactly \(R_U\).
For surjectivity, take a tuple with common value \(d\in D\)
and choose one \(g\in\Gamma\) over \(d\). Multiplying by
the inverse of its actual tuple leaves an arbitrary element of
\(\prod_pU^{(p)}\); \textup{(E.58.8)} supplies its source
preimage. This proves \textup{(E.58.2)} without choosing a
splitting of any of its extensions.

In particular, local automorphisms inducing one common base
automorphism, together with local peripheral conjugators
having common \(D\)-values, define a peripheral automorphism
of the displayed fibre product. The relative lifting theorem
retains each of these full local maps and conjugators.
This observation will be used with the identity automorphism
in every component.

\paragraph{The actual finite Fermat source.}
The subgroup \(F\) is characteristic in \(\mathcal B\).
Indeed the six epimorphisms \(C_2*C_3\to S_3\) send the
two generators to a transposition and a three-cycle; they
form one orbit under \(\operatorname{Aut}(S_3)\) and have
one kernel. The same assertion holds for the continuous
maps from the profinite completion. Hence \(H=F'F^3\)
is characteristic as well. Its index is
\[
 [\mathcal B:H]=6\cdot3^2=54,
 \qquad F/H\simeq(\mathbf Z/3)^2.
 \tag{E.58.9}
\]

Let \(V_4\) be the four matrices modulo four
\[
 I,\qquad
 \begin{pmatrix}1&2\\0&1\end{pmatrix},\qquad
 \begin{pmatrix}1&0\\2&1\end{pmatrix},\qquad
 \begin{pmatrix}1&2\\2&1\end{pmatrix}.
 \tag{E.58.10}
\]
Every element of \(q(F)\) has a unique representative modulo
twelve whose reduction modulo four belongs to \(V_4\).
Multiplication of those representatives gives the actual
isomorphism
\[
 q(F)\ \xrightarrow{\ \sim\ }\ 
 V_4\times\operatorname{SL}_2(\mathbf F_3),
 \qquad [M]\longmapsto(M\bmod4,M\bmod3).
 \tag{E.58.11}
\]
The uniqueness follows from \(-I\notin V_4\), and the
Chinese remainder map proves injectivity. The literal
\(x,y\) matrices generate all 96 elements of this product,
as verified independently in the certificate below; this
also identifies \(q(F)\) with the kernel of reduction to
\(\operatorname{SL}_2(\mathbf F_2)=S_3\).

The derived subgroup of \(\operatorname{SL}_2(\mathbf F_3)\)
is its quaternion subgroup \(Q_8\). One explicit pair of
quaternion generators is
\[
 i=\begin{pmatrix}0&1\\2&0\end{pmatrix},\qquad
 j=\begin{pmatrix}1&1\\1&2\end{pmatrix};\qquad
 i^2=j^2=-I,\quad ij=-ji.
 \tag{E.58.12}
\]
Their group has order eight and quotient of order three;
the literal commutator and its conjugates generate it.
Cubes kill that quotient, whereas cubing is the identity
on \(V_4\). Thus \textup{(E.58.11)} gives exactly
\[
                   q(H)=V_4\times Q_8,
                   \qquad |q(H)|=32.
 \tag{E.58.13}
\]
In particular the map \(q\) kills \(O^2(H)\), so it
factors through the actual source
\(\Gamma_2=\mathcal B/O^2(H)\).

Set \(D=\mathcal B/H\) and retain the actual joint image
\[
 \begin{gathered}
 J=\operatorname{im}(\mathcal B\to Q\times D)
       =Q\times_{Q/q(H)}D,\\
 |Q/q(H)|=18,\qquad |J|=32\cdot54=1728.
 \end{gathered}
 \tag{E.58.14}
\]
The map \(D\to Q/q(H)\) is induced by \(q\).
To prove the fibre-product equality, lift the prescribed
\(D\)-coordinate to \(\mathcal B\); the discrepancy in
its \(Q\)-coordinate lies in \(q(H)\), and an actual
element of \(H\) corrects it. The kernel over \(D\)
is consequently exactly \(q(H)\).
This joint source also factors through \(\Gamma_2\),
and the kernel of
\[
                         \Gamma_2\twoheadrightarrow J
 \tag{E.58.15}
\]
is a closed subgroup of \(H^{(2)}\), hence is pro-two.

For an independent explicit model of \(D\), use affine
transformations of \(\mathbf F_3^2\), written \((L,v)\),
with multiplication \((L,v)(M,w)=(LM,v+Lw)\).
The images
\[
 s\longmapsto
   \left(\begin{pmatrix}0&1\\1&0\end{pmatrix},\binom00\right),
 \qquad
 r\longmapsto
   \left(\begin{pmatrix}0&2\\1&2\end{pmatrix},\binom01\right)
 \tag{E.58.16}
\]
satisfy the defining torsion relations. The literal
\(x,y\) become translations by the two standard basis
vectors. Their image has order 54, so its kernel is
exactly \(H\), by \textup{(E.58.9)}. The finite
certificate enumerates the joint source through these
same generator maps.

\paragraph{A specified relative centralizer element.}
The following computation is in the fixed matrices of
\textup{(E.58.3)}:
\[
 \begin{gathered}
 q(x)=\left[\begin{pmatrix}1&2\\0&1\end{pmatrix}\right],\qquad
 q(y)=\left[\begin{pmatrix}1&0\\-2&1\end{pmatrix}\right],\\
 q([x,y])=\left[\begin{pmatrix}1&8\\8&5\end{pmatrix}\right],\qquad
 z_0:=q([x,y]^2)=[5I].
 \end{gathered}
 \tag{E.58.17}
\]
Thus \(z_0\) is central of order two in \(Q\), and
\(z_0\in q(H)\) has coordinates \((I,-I)\) under
\textup{(E.58.11)}. The literal word \([x,y]^2\in H\)
shows that
\[
                         (z_0,1)\in C_J(r).
 \tag{E.58.18}
\]

For clarity, here is the precise centralizer lifting
argument. In a surjection \(G\to E\) with pro-\(p\)
kernel \(K\), let \(a\in G\) have finite order prime
to \(p\). Its image has the same order. Given any
\(e\in C_E(\bar a)\), choose a lift \(g\).
The subgroups \(\langle a\rangle\) and
\(\langle gag^{-1}\rangle\) are complements to \(K\)
in the preimage of \(\langle\bar a\rangle\).
At each finite quotient, coprime complement conjugacy
gives an element of the kernel conjugating them with
the same generator image. The resulting closed
conjugator sets in the compact group \(K\) are nonempty
and have the finite intersection property. Their
intersection gives \(u\in K\) with
\(ugag^{-1}u^{-1}=a\). Thus \(ug\) centralizes
\(a\) and lifts the prescribed \(e\).

Apply this argument to \textup{(E.58.15)}, \(p=2\),
\(a=r\), and \(e=(z_0,1)\). The element \(r\) has
order three in both groups, because the retained
\(S_3\) quotient detects that order. We obtain an
actual element
\[
 c\in C_{\Gamma_2}(r)\cap H^{(2)},\qquad
 c\longmapsto(z_0,1)\text{ in }J.
 \tag{E.58.19}
\]
More explicitly, it has the form \(u[x,y]^2\), with
\(u\in\ker(\Gamma_2\to J)\) supplied by the preceding
compact conjugator construction. This specifies its
actual maps, rather than replacing a finite centralizer
element by an unproved full centralizer lift.

\paragraph{The full gentle automorphism and its exact frames.}
Use \textup{(E.58.2)} with this \(U=H\). Take the
identity automorphism in every local component. Choose
the peripheral conjugators
\[
 b_s=1,\qquad b_{\tau}=1,\qquad
 b_r=(c\text{ at }2,\ 1\text{ at every other prime}).
 \tag{E.58.20}
\]
All components of \(b_r\) have value one in \(D\),
so this is an actual element of the indicated fibre
product. It centralizes \(r\) in every component.
Consequently these are literal peripheral conjugators
for the identity automorphism of
\(\mathcal B/R_{\rm nil}(H)\).

The relative lifting theorem gives an automorphism of
\(\mathcal B\) with these full relative coordinates.
Conjugate it by the inverse of its \(\tau\)-conjugator,
which lies in \(R_{\rm nil}(H)\), to obtain \(\Phi\)
fixing \(\tau\). Its \(s\)- and \(r\)-conjugators
\(\beta,A\) still have precisely the prescribed
images \(1,b_r\). Their values in \(D\) are one,
so they are actual elements of \(H\subset F\).
Since \(q\) factors through the two-primary component,
their finite values are
\(q(\beta)=1\), \(q(A)=z_0\), as in
\textup{(E.58.5)}.

Put \(f=s\Phi(s)\). The retained relative automorphism
is the identity, so \(f\in R_{\rm nil}(H)\subset H'\subset F'\).
Exactly as in D.44, the torsion and product relations give
\[
 \begin{gathered}
 \Phi(s)=sf,\qquad \Phi(r)=f^{-1}r,\qquad
 \Phi(x)=x,\qquad \Phi(y)=f^{-1}yf,\\
 f(y,x)f(x,y)=1,\qquad
 f((xy)^{-1},x)f(y,(xy)^{-1})f(x,y)=1.
 \end{gathered}
 \tag{E.58.21}
\]
The restriction to the characteristic subgroup \(F\)
is an actual automorphism. These are the multiplier-one
full gentle equations, so \((1,f)\in\widehat{\mathrm{GT}}_0\).

The full accessory frames here have no hidden pure
centralizer ambiguity. Herfort--Ribes, Theorem B'
(Section 4, printed page 159), proves for a free
profinite product that the centralizer of a nonidentity
factor element is its centralizer in that factor.
Applied to the two specified factors of \(\mathcal B\),
this gives
\[
 C_{\mathcal B}(s)=\langle s\rangle,\qquad
 C_{\mathcal B}(r)=\langle r\rangle,
 \qquad C_{\mathcal B}(s)\cap F=C_{\mathcal B}(r)\cap F=1.
 \tag{E.58.22}
\]
The last equalities use the injective images of those
factors in the specified \(S_3\) quotient.
Thus \(\beta,A\in F\) are the unique pure conjugators
for this normalized full automorphism. In particular
the nontrivial relative \(r\)-frame is retained by
an actual full automorphism; it has not been identified
with a full centralizer element.\cite{BCHerfortRibes1985}

\paragraph{The positive CM convention and the resulting coset.}
Use exactly E.1's representatives and coordinates
\(T=A^{-1}V\), \(W=\beta^{-1}V\). Here \(V=1\), so
\[
 q(T)=[5I],\qquad q(W)=1.
 \tag{E.58.23}
\]
The unique representative of \([5I]\) whose mod-four
value belongs to \(V_4\) is \(5I\) itself.
Its inverse is again \(5I\), with mod-three value
\(-I\). For the specified positive matrix
\(R=\left(\begin{smallmatrix}0&-1\\1&-1\end{smallmatrix}\right)\),
the rational three-torsion line \(\ker(R-I)\) modulo
three is generated by \((1,2)^t\). The \(A^{-1}\)
coordinate sends this vector to \((2,1)^t\), giving
E.1's first bit one. The \(\beta^{-1}\) coordinate
is the identity, giving its second bit zero.
Thus the two restrictions compared in E.1 disagree.
Through the precise finite-coordinate comparison of
E.1 and LM7 D.25, this is the extra CM coset in Task E.2.

The finite centralizer calculation is essential for
these coordinates: \(\Phi\) induces the identity on
\(Q\), while its unique full \(r\)-accessory frame
has value \([5I]\). The preceding construction and
\textup{(E.58.22)} prove exactly how those two statements
hold for the same actual automorphism.

\paragraph{Comparison with the original homology shadows.}
For this construction one does not need a factorization
of the entire \(G_6\) source through \(\Gamma_2\).
In fact no such factorization exists for the specified
maps. The element
\[
              x^6=\tau^{-12}\in H\cap N_B
 \tag{E.58.24}
\]
has primitive \(N_B\)-abelianization. To check this
integrally, use the Bass--Serre graph of E.1. The chain
of \(\tau^{12}\) is
\[
       \sum_{j=0}^{11}
          \bigl(e_{q(\tau)^j q(s)}-e_{q(\tau)^{j+1}}\bigr).
 \tag{E.58.25}
\]
It has zero boundary and 24 distinct nonzero edge
coefficients, each \(1\) or \(-1\). Hence it is
primitive in the integral cycle lattice. Its image
in \(G_6=B/(N_B'N_B^6)\) therefore has order six.
But the image of \(H\) in \(\Gamma_2\) is pro-two,
so every finite image of its element has two-power
order. This proves the asserted failure of that
factorization through an explicit element.

The full automorphism just constructed nevertheless
has an actual \(G_6\) shadow. Indeed \(q\Phi=q\),
so \(\Phi\) preserves \(N_B\) and its characteristic
subgroup \(N_B'N_B^6\). It therefore induces an actual
automorphism of this specified \(G_6\), and that shadow
has the finite frames
\textup{(E.58.5)}. E.9 excludes the particular earlier
recorded representatives; the present construction
gives a full-source point in the whole CM matrix
fibre. It does not identify that point with one of
those excluded representatives.

The independent standard-library program
\texttt{bc\_relative\_cm\_sources.py} enumerates the
576-element projective matrix source, its exact
96- and 32-element subgroup images, the affine
54-element quotient, and the full 1728-element
joint source. It checks every product in
\textup{(E.58.11)}, the literal word
\([x,y]^2\), the positive signs \((1,0)\), and
the primitive integral chain in
\textup{(E.58.25)}. The infinite centralizer and
peripheral lifts are the complete proofs above.
\end{proof}
